\documentclass[a4paper,14pt]{extarticle}

% Компилировать XeLaTeX или LuaLaTeX



\usepackage[T2A]{fontenc}
\usepackage[utf8]{inputenc}
\usepackage[russian]{babel}
\usepackage{amsmath, amssymb}
\usepackage{geometry}
\usepackage{graphicx}
\usepackage{float}
\usepackage{booktabs}
\usepackage{array}
\usepackage{hyperref}
\usepackage{listings}
\usepackage{xcolor}
\usepackage{caption}
\usepackage{listings}
\usepackage{xcolor}

\definecolor{codegreen}{rgb}{0,0.6,0}
\definecolor{codegray}{rgb}{0.5,0.5,0.5}
\definecolor{codepurple}{rgb}{0.58,0,0.82}
\definecolor{backcolour}{rgb}{0.96,0.96,0.94}
\lstdefinestyle{pythonstyle}{
    backgroundcolor=\color{backcolour},   
    commentstyle=\color{codegreen}\itshape,
    keywordstyle=\color{blue}\bfseries,
    numberstyle=\tiny\color{codegray},
    stringstyle=\color{codepurple},
    basicstyle=\ttfamily\small,
    breakatwhitespace=false,         
    breaklines=true,                 
    captionpos=b,                    
    keepspaces=true,                 
    numbers=left,                    
    numbersep=5pt,                  
    showspaces=false,                
    showstringspaces=false,
    showtabs=false,                  
    tabsize=4,
    language=Python,
    % Поддержка кириллицы в комментариях и строках:
    extendedchars=true,
    keepspaces=true
}

\geometry{left=30mm,right=15mm,top=20mm,bottom=20mm}
\setlength{\parindent}{1.25cm}
\setlength{\parskip}{0.3em}
\renewcommand{\baselinestretch}{1.15}

\lstset{
    basicstyle=\ttfamily\small,
    breaklines=true,
    frame=single,
    numbers=left,
    numberstyle=\tiny,
    keywordstyle=\color{blue},
    commentstyle=\color{gray},
    stringstyle=\color{teal},
    showstringspaces=false,
    tabsize=4
}

\newcommand{\placeholder}[1]{%
\begin{center}
\fbox{%
\begin{minipage}[c][0.18\textheight][c]{0.88\textwidth}
\centering #1
\end{minipage}}
\end{center}}

\begin{document}

\begin{titlepage}
    \begin{center}
        Министерство науки и высшего образования Российской Федерации\\
        Университет ИТМО\\[2cm]

        \textbf{Лабораторная работа №4}\\[0.3cm]
        \textbf{«Аппроксимация функции методом наименьших квадратов»}\\[0.5cm]
        Вариант 3\\[3cm]
    \end{center}

    \begin{flushright}
        Выполнил: Григорьев Даниил Александрович\\
        Группа: P3216\\
        Проверил: Наумова Надежда Александровна\\
    \end{flushright}

    \vfill
    \begin{center}
        Санкт-Петербург\\
        2026
    \end{center}
\end{titlepage}

\tableofcontents
\newpage

\section{Цель работы}

Цель работы: найти функцию, являющуюся наилучшим приближением заданной табличной функции по методу наименьших квадратов, сравнить несколько видов аппроксимирующих зависимостей и выбрать лучшую по среднеквадратическому отклонению.

\section{Задание}

Для вычислительной части используется вариант 3:
\[
    y = \frac{4x}{x^4 + 3}, \qquad x \in [-2; 0], \qquad h = 0.2.
\]

Необходимо сформировать таблицу значений функции на заданном интервале, построить линейное и квадратичное приближения по 11 точкам, найти среднеквадратические отклонения и выбрать лучшее приближение.

В программной части необходимо исследовать следующие зависимости:
\begin{itemize}
    \item линейная функция;
    \item полиномиальная функция 2-й степени;
    \item полиномиальная функция 3-й степени;
    \item экспоненциальная функция;
    \item логарифмическая функция;
    \item степенная функция.
\end{itemize}

Программа должна принимать таблицу из 8--12 точек, строить графики функций, выводить коэффициенты, меру отклонения, среднеквадратическое отклонение, массивы $x_i$, $y_i$, $\varphi(x_i)$, $\varepsilon_i$, коэффициент корреляции Пирсона для линейной зависимости и коэффициент детерминации $R^2$.

\section{Рабочие формулы метода}

Метод наименьших квадратов выбирает коэффициенты аппроксимирующей функции так, чтобы сумма квадратов отклонений была минимальной:
\[
    S = \sum_{i=1}^{n} \left(\varphi(x_i) - y_i\right)^2 \to \min.
\]

Среднеквадратическое отклонение вычисляется по формуле:
\[
    \sigma = \sqrt{\frac{S}{n}}.
\]

Для линейной зависимости $\varphi(x) = ax + b$ система нормальных уравнений имеет вид:
\[
\begin{cases}
    a\sum x_i^2 + b\sum x_i = \sum x_i y_i,\\
    a\sum x_i + bn = \sum y_i.
\end{cases}
\]

Для полинома степени $m$ используется вид:
\[
    \varphi(x) = a_0 + a_1x + a_2x^2 + \dots + a_mx^m.
\]

Коэффициенты находятся из системы нормальных уравнений:
\[
    \sum_{j=0}^{m} a_j \sum_{i=1}^{n} x_i^{j+k} = \sum_{i=1}^{n} y_i x_i^k,
    \qquad k = 0, 1, \dots, m.
\]

Неполиномиальные зависимости приводятся к линейному виду:
\begin{itemize}
    \item экспоненциальная: $y = ae^{bx}$, тогда $\ln y = \ln a + bx$;
    \item логарифмическая: $y = a\ln x + b$;
    \item степенная: $y = ax^b$, тогда $\ln y = \ln a + b\ln x$.
\end{itemize}

Коэффициент детерминации:
\[
    R^2 = 1 - \frac{\sum_{i=1}^{n}(y_i - \varphi(x_i))^2}{\sum_{i=1}^{n}(y_i - \overline{y})^2}.
\]

Коэффициент корреляции Пирсона для линейной зависимости:
\[
    r = \frac{\sum_{i=1}^{n}(x_i - \overline{x})(y_i - \overline{y})}
    {\sqrt{\sum_{i=1}^{n}(x_i - \overline{x})^2 \sum_{i=1}^{n}(y_i - \overline{y})^2}}.
\]

\section{Вычислительная часть}

Составил таблицу из 11 точек для функции

\begin{table}[H]
\centering
\caption{Таблица значений функции $y = \frac{4x}{x^4 + 3}$}
\begin{tabular}{c|c|c}
\toprule
$i$ & $x_i$ & $y_i$ \\
\midrule
1 & -2.0 & -0.421 \\
2 & -1.8 & -0.533 \\
3 & -1.6 & -0.670 \\
4 & -1.4 & -0.819 \\
5 & -1.2 & -0.946 \\
6 & -1.0 & -1.000 \\
7 & -0.8 & -0.939 \\
8 & -0.6 & -0.767 \\
9 & -0.4 & -0.529 \\
10 & -0.2 & -0.267 \\
11 & 0.0 & 0.000 \\
\bottomrule
\end{tabular}
\end{table}

Линейное приближение имеет вид:
\[
    \varphi_1(x) = 0.168499x - 0.457840.
\]

Для линейного приближения:
\[
    S_1 \approx 0.855023, \qquad \sigma_1 \approx 0.279.
\]

Квадратичное приближение имеет вид:
\[
    \varphi_2(x) = 0.773805x^2 + 1.716109x + 0.006443.
\]

Для квадратичного приближения:
\[
    S_2 \approx 0.033026, \qquad \sigma_2 \approx 0.055.
\]

Так как $\sigma_2 < \sigma_1$, в вычислительной части лучшим из двух приближений является квадратичное.

\section{Блок-схемы реализованного алгоритма}


\begin{figure}[H]
    \centering
    \includegraphics[width=0.75\linewidth]{1.png}
    \caption{Линейная аппроксимация}
    \label{fig:placeholder}
\end{figure}


\begin{figure}[H]
    \centering
    \includegraphics[width=0.75\linewidth]{2.png}
    \caption{Полиномиальная аппроксимация 2 степени}
    \label{fig:placeholder}
\end{figure}

\begin{figure}[H]
    \centering
    \includegraphics[width=0.75\linewidth]{3.png}
    \caption{Полиномиальная аппроксимация 3 степени}
    \label{fig:placeholder}
\end{figure}
\begin{figure}[H]
    \centering
    \includegraphics[width=0.75\linewidth]{4.png}
    \caption{Экспоненциальная аппроксимация}
    \label{fig:placeholder}
\end{figure}


\begin{figure}[H]
    \centering
    \includegraphics[width=0.75\linewidth]{5.png}
    \caption{Логарифмическая аппроксимация}
    \label{fig:placeholder}
\end{figure}


\begin{figure}[H]
    \centering
    \includegraphics[width=0.75\linewidth]{6.png}
    \caption{Степенная аппроксимация}
    \label{fig:placeholder}
\end{figure}



\section{Ссылка на GitHub с основной реализацией}

\noindent Ссылка на репозиторий: 
\url{https://github.com/Megadeth2006/Lab4_MathComp}


\section{Пример работы программы}
\begin{figure}[H]
    \centering
    \includegraphics[width=0.75\linewidth]{screen2.png}
    \caption{Главное окно после запуска}
    \label{fig:placeholder}
\end{figure}
\begin{figure}[H]
    \centering
    \includegraphics[width=0.75\linewidth]{screen1.png}
    \caption{Отображение апппроксимации значений из варианта}
    \label{fig:placeholder}
\end{figure}

Для набора данных варианта 3 программная часть строит все применимые функции. Экспоненциальная, логарифмическая и степенная зависимости для данного набора не применяются, так как в таблице есть неположительные значения $x$ и $y$. Среди применимых моделей программа выбирает лучшую по минимальному среднеквадратическому отклонению.

\section{Тестирование программы}

Программа была проверена на следующих наборах данных:
\begin{enumerate}
    \item корректный набор варианта 3 из 11 точек;
    \item положительный набор данных, для которого применимы все 6 видов аппроксимации;
    \item некорректный набор данных: недостаточное количество точек и неверный числовой ввод.
\end{enumerate}

\begin{figure}[H]
    \centering
    \includegraphics[width=0.75\linewidth]{screen3.png}
    \caption{Корректный набор из варианта 3}
    \label{fig:placeholder}
\end{figure}
\begin{figure}[H]
    \centering
    \includegraphics[width=0.75\linewidth]{screen4.png}
    \caption{положительный набор данных.}
    \label{fig:placeholder}
\end{figure}

\begin{figure}[H]
    \centering
    \includegraphics[width=0.75\linewidth]{screen5.png}
    \caption{человекочитаемое сообщение об ошибке при некорректном вводе}
    \label{fig:placeholder}
\end{figure}


\section{Листинг}
\begin{lstlisting}[style=pythonstyle, caption={Модуль вычислительных методов аппроксимации (approximation\_methods.py)}, label={lst:approximation_methods}]
from __future__ import annotations
from typing import Callable, Any
import math
from dataclasses import dataclass


@dataclass(frozen=True)
class FittedFunction:
    expression: str
    coefficients: dict[str, float]
    evaluator: Callable[..., Any]


EPSILON = 1e-12


def solve_linear_system(matrix: list[list[float]], vector: list[float]) -> list[float]:
    size = len(vector)
    a = [row[:] for row in matrix]
    b = vector[:]

    for column in range(size):
        pivot_row = max(range(column, size), key=lambda row: abs(a[row][column]))
        pivot = a[pivot_row][column]
        if abs(pivot) < EPSILON:
            raise ValueError("Система нормальных уравнений вырождена. Проверьте исходные точки.")

        if pivot_row != column:
            a[column], a[pivot_row] = a[pivot_row], a[column]
            b[column], b[pivot_row] = b[pivot_row], b[column]

        pivot = a[column][column]
        for j in range(column, size):
            a[column][j] /= pivot
        b[column] /= pivot

        for row in range(size):
            if row == column:
                continue
            factor = a[row][column]
            if abs(factor) < EPSILON:
                continue
            for j in range(column, size):
                a[row][j] -= factor * a[column][j]
            b[row] -= factor * b[column]

    return b


def least_squares_by_basis(xs: list[float], ys: list[float], basis_functions: list[Callable[..., Any]]) -> list[float]:
    size = len(basis_functions)
    matrix: list[list[float]] = []
    vector: list[float] = []

    basis_values = [[basis(x) for x in xs] for basis in basis_functions]

    for i in range(size):
        matrix_row: list[float] = []
        for j in range(size):
            matrix_row.append(sum(basis_values[i][k] * basis_values[j][k] for k in range(len(xs))))
        matrix.append(matrix_row)
        vector.append(sum(ys[k] * basis_values[i][k] for k in range(len(xs))))

    return solve_linear_system(matrix, vector)


def fit_polynomial(xs: list[float], ys: list[float], degree: int) -> FittedFunction:
    basis_functions = [lambda x, power=power: x ** power for power in range(degree + 1)]
    coefficients = least_squares_by_basis(xs, ys, basis_functions)

    def evaluate(x: float) -> float:
        return sum(coefficients[power] * x ** power for power in range(degree + 1))

    names = ["a0", "a1", "a2", "a3"]
    coefficient_map = {names[index]: coefficients[index] for index in range(degree + 1)}

    if degree == 1:
        expression = "a1*x + a0"
    elif degree == 2:
        expression = "a2*x^2 + a1*x + a0"
    elif degree == 3:
        expression = "a3*x^3 + a2*x^2 + a1*x + a0"
    else:
        expression = "polynomial"

    return FittedFunction(expression, coefficient_map, evaluate)


def fit_exponential(xs: list[float], ys: list[float]) -> FittedFunction:
    if any(y <= 0 for y in ys):
        raise ValueError("экспоненциальная аппроксимация y = a*e^(b*x) требует y > 0 для всех точек")

    transformed_y = [math.log(y) for y in ys]
    c0, c1 = least_squares_by_basis(xs, transformed_y, [lambda _x: 1.0, lambda x: x])
    a = math.exp(c0)
    b = c1

    def evaluate(x: float) -> float:
        return a * math.exp(b * x)

    return FittedFunction("a*e^(b*x)", {"a": a, "b": b}, evaluate)


def fit_logarithmic(xs: list[float], ys: list[float]) -> FittedFunction:
    if any(x <= 0 for x in xs):
        raise ValueError("логарифмическая аппроксимация y = a*ln(x) + b требует x > 0 для всех точек")

    c0, c1 = least_squares_by_basis(xs, ys, [lambda _x: 1.0, lambda x: math.log(x)])
    b = c0
    a = c1

    def evaluate(x: float) -> float:
        return a * math.log(x) + b

    return FittedFunction("a*ln(x) + b", {"a": a, "b": b}, evaluate)


def fit_power(xs: list[float], ys: list[float]) -> FittedFunction:
    if any(x <= 0 for x in xs):
        raise ValueError("степенная аппроксимация y = a*x^b требует x > 0 для всех точек")
    if any(y <= 0 for y in ys):
        raise ValueError("степенная аппроксимация y = a*x^b требует y > 0 для всех точек")

    transformed_x = [math.log(x) for x in xs]
    transformed_y = [math.log(y) for y in ys]
    c0, c1 = least_squares_by_basis(transformed_x, transformed_y, [lambda _x: 1.0, lambda x: x])
    a = math.exp(c0)
    b = c1

    def evaluate(x: float) -> float:
        return a * x ** b

    return FittedFunction("a*x^b", {"a": a, "b": b}, evaluate)
\end{lstlisting}
\section{Вывод}

В ходе лабораторной работы был реализован метод наименьших квадратов для аппроксимации табличной функции. Программа исследует линейную, квадратичную, кубическую, экспоненциальную, логарифмическую и степенную зависимости, вычисляет меру отклонения, среднеквадратическое отклонение, коэффициент детерминации и коэффициент корреляции Пирсона для линейной функции.

Для вычислительной части варианта 3 квадратичное приближение оказалось лучше линейного, так как имеет меньшее среднеквадратическое отклонение. В программной части лучшая функция выбирается автоматически среди всех применимых моделей.

\end{document}
